\chapter{Éléments complémentaires}

\section{Structure du dépôt}

\begin{lstlisting}[style=mystyle,caption={Organisation du dépôt de code},captionpos=b]
sovereign-legal-rag/
|
|-- data/
|   |-- raw/                        # PDF sources, non modifies
|   |-- processed/
|   |   `-- corpus_chunks.jsonl     # Corpus structure, un article par ligne
|   `-- chroma/                     # Base vectorielle (generee)
|
|-- src/
|   |-- parser.py                   # PDF -> corpus JSONL
|   |-- database.py                 # Corpus -> index vectoriel
|   |-- retriever.py                # Recherche hybride (dense + BM25 + RRF)
|   |-- generator.py                # Generation ancree via Ollama
|   `-- app.py                      # Orchestrateur : abstention + citations
|
|-- eval/
|   |-- reference_qa.jsonl          # Jeu de reference (37 questions)
|   |-- run_eval.py                 # Recall@k, MRR (sans LLM)
|   |-- run_answer_eval.py          # Criteres de succes
|   |-- run_measures.py             # Latence, taxonomie, figures
|   `-- figures/                    # Figures SVG generees
|
|-- web/                            # Interface React
|-- serve.py                        # Serveur web local
|-- demo.py                         # Demonstration scriptee
`-- JOURNAL.md                      # Carnet de bord quotidien
\end{lstlisting}

\section{Reproduction des résultats}

L'ensemble des chiffres de ce rapport se régénère par les commandes suivantes, exécutées depuis
la racine du dépôt :

\begin{lstlisting}[language=bash,style=mystyle,caption={Commandes de reproduction},captionpos=b]
# 1. Construction du corpus a partir du PDF officiel
python src/parser.py

# 2. Indexation (telecharge le modele d'embeddings au premier lancement)
python src/database.py

# 3. Metriques de recherche (ne necessite pas le LLM)
python eval/run_eval.py

# 4. Evaluation des reponses generees (necessite Ollama)
python eval/run_answer_eval.py

# 5. Mesures de latence et figures
python eval/run_measures.py
\end{lstlisting}

\section{Le prompt système}

Le prompt qui encadre la génération constitue l'un des éléments les plus déterminants du système.
Sa version finale impose les règles suivantes :

\begin{enumerate}
    \item Répondre uniquement à partir des articles fournis, sans recourir à aucune autre
    connaissance.
    \item Citer systématiquement les numéros d'articles utilisés.
    \item Employer une phrase d'abstention exacte lorsque les articles ne permettent pas de
    répondre.
    \item Ne jamais compléter un aveu d'absence par une supposition sur le contenu d'un article
    manquant.
    \item Répondre par défaut lorsque les articles fournis contiennent de quoi le faire ---
    l'abstention étant réservée aux cas où le sujet n'est réellement pas traité.
    \item Ne jamais citer un numéro de loi ou d'article provenant d'un autre code.
    \item Ne jamais donner de conseil juridique personnalisé.
\end{enumerate}

Les règles 4, 5 et 6 ont été ajoutées en réaction à des défaillances observées lors de
l'évaluation, décrites en section~\ref{sec:hallucination}.

% TODO : ajouter ici, si utile, des extraits de code commentés
% (parser, fusion RRF, verification des citations) ou des captures
% de l'interface web.
